\documentclass[11pt,a4paper]{article}

\usepackage[T1]{fontenc}
\usepackage[utf8]{inputenc}
\usepackage{lmodern}
\usepackage{geometry}
\usepackage{booktabs}
\usepackage{listings}
\usepackage{xcolor}
\usepackage{hyperref}
\usepackage{amsmath}

\geometry{margin=2.4cm}
\hypersetup{colorlinks=true, linkcolor=black, urlcolor=blue}

\lstset{
  basicstyle=\ttfamily\small,
  breaklines=true,
  frame=single,
  backgroundcolor=\color{gray!6},
  rulecolor=\color{gray!35},
  showstringspaces=false,
  columns=fullflexible
}

\title{Short Project Report\\Differential-Drive Robot with ESP32-C6 and Future ESP32-S3 Camera Extension}
\author{Krzysztof Sawicki, Michal Szerechowicz, Mateusz Skorupinski}
\date{\today}

\begin{document}
\maketitle

\section{Project Goal}

The goal of the project was to build a compact mobile robot that can be
controlled wirelessly, move as a differential-drive platform, keep its heading
stable using an inertial sensor, and provide a simple browser interface for
manual and semi-autonomous operation. The current robot is based on the
ESP32-C6, placed on the CodeCell C6 Drive board, which is responsible for
driving the robot and running the control logic. An ESP32-S3 camera module was
also prepared in software, but it is not currently mounted or used on the
robot. It is intended as a future extension that can provide a Wi-Fi access
point, live video stream, snapshot endpoint, and camera power/stream controls.

\section{Mechanical Design}

The robot uses a differential-drive layout. Two DC motors are mounted on the
left and right sides of a 3D-printed chassis. By changing the relative speed
and direction of these motors the robot can drive forward, reverse, follow a
curve, or rotate in place. The chassis files are included in the project as STL
models for both a differential-drive and a four-wheel variant. The implemented
control software is written for the two-motor differential-drive arrangement.

The main mechanical requirements were compactness, repeatability, and short
wiring. The CodeCell C6 Drive board and LiPo battery are mounted centrally to
keep the mass close to the robot's centre. This helps both stability and
heading control, because the robot rotates more predictably when the mass is
not strongly biased to one side.

\section{Electronics}

The drive robot is built around the CodeCell C6 Drive board. It integrates the
ESP32-C6 microcontroller, wireless communication, a BNO085 9-axis IMU, light
and proximity sensing, a LiPo/USB-C power path, and dual H-bridge motor drivers.
This reduced the amount of custom wiring because the same board handles
control, sensing, communication, battery readings, and motor output.

\begin{table}[h]
\centering
\begin{tabular}{@{}ll@{}}
\toprule
Module & Role in the project \\
\midrule
ESP32-C6 / CodeCell C6 Drive & Main robot controller and motor driver \\
BNO085 IMU & Yaw/heading measurement for stabilised movement \\
Dual H-bridge drivers & Independent control of left and right DC motors \\
LiPo power path & Battery supply and voltage/level reporting \\
ESP32-S3 camera board & Future camera extension; not currently used on robot \\
\bottomrule
\end{tabular}
\end{table}

\section{ESP32-C6 Drive Firmware}

The C6 firmware is an Arduino/PlatformIO project. It uses the CodeCell library
and separates the robot behaviour into three FreeRTOS tasks:

\begin{itemize}
  \item \textbf{MotorTask} runs the low-level H-bridge output and applies the
  latest signed motor commands to Drive1 and Drive2.
  \item \textbf{ControlTask} reads the active command, runs the motion logic,
  updates the IMU, computes motor output, and publishes robot status.
  \item \textbf{CommunicationTask} starts the robot's Wi-Fi access point and
  serves the REST API used by the browser controller.
\end{itemize}

The tasks communicate through shared command, motor, and status buffers. Each
buffer is protected by a mutex, which avoids inconsistent reads when one task
updates a structure while another task is using it.

\begin{lstlisting}[language=C++,caption={Task startup on the ESP32-C6}]
xTaskCreate(motorTask, "MotorTask", 2048, NULL, 3, NULL);
xTaskCreate(controlTask, "ControlTask", 4096, NULL, 2, NULL);
xTaskCreate(communicationTask, "CommunicationTask", 8192, NULL, 1, NULL);
\end{lstlisting}

\section{Motor Mixing and Movement}

Manual motion is based on two normalized commands: \texttt{move} and
\texttt{rotation}, both limited to the range from -100 to 100. The controller
mixes them into left and right motor values:

\[
m_1 = \operatorname{clip}(\texttt{move} + \texttt{rotation}), \qquad
m_2 = \operatorname{clip}(\texttt{move} - \texttt{rotation})
\]

The sign of each result selects direction and the absolute value selects motor
speed. Equal values drive the robot straight. A difference between the motors
causes a turn, and opposite signs rotate the robot in place. In software this
is implemented by \texttt{calculateMotorMix()}, with final output stored as a
direction bit and an unsigned PWM-like speed value.

To improve reliability at low speed, the firmware adds a start ramp. Small DC
motors often need a short torque boost to overcome static friction. The robot
therefore begins with a stronger command and smoothly decays toward the
requested speed during the first 500 ms of motion.

\section{Heading Stabilisation}

The BNO085 IMU provides the robot's yaw angle. The controller compares the
measured yaw with a target yaw and uses a PID controller to compute a rotation
correction. The angular error is wrapped to the shortest path, so an error
near the -180/180 degree boundary is handled correctly.

The active tuning constants in the C6 firmware are:

\begin{center}
\begin{tabular}{@{}ll@{}}
\toprule
Parameter & Value \\
\midrule
$K_p$ & 0.25 \\
$K_i$ & 0.02 \\
$K_d$ & 0.03 \\
Integral limit & $\pm 120$ \\
Maximum correction & $\pm 45$ \\
\bottomrule
\end{tabular}
\end{center}

\begin{lstlisting}[language=C++,caption={Core yaw correction idea}]
float err = angleError(targetYaw, yaw);
yawIntegral = constrain(yawIntegral + err * dt,
                        -YAW_INTEGRAL_LIMIT,
                         YAW_INTEGRAL_LIMIT);
float correction = YAW_KP * err
                 + YAW_KI * yawIntegral
                 + YAW_KD * derivative;
\end{lstlisting}

This heading loop is used in both direct and target modes. In direct mode the
joystick rotation changes the yaw setpoint over time, so the robot turns
relative to its current direction while still using IMU feedback. In target
mode the heading loop keeps the robot aligned during autonomous travel.

\section{Control Modes}

The robot supports two operating modes.

\textbf{Direct mode} is used for real-time driving from the browser joystick.
The browser periodically sends \texttt{/move?move=...\&rotation=...}. If these
commands stop arriving for more than 300 ms, the C6 firmware stops the motors.
This timeout prevents uncontrolled motion after Wi-Fi or browser failure.

\textbf{Target mode} accepts a distance, relative angle, and optional IMU flag
through the \texttt{/target} endpoint. The robot first rotates toward the
requested heading. When the yaw error is small enough, it drives forward or
backward while the PID loop maintains the heading. Distance can be estimated by
a calibrated time model or by integrating the linear acceleration from the IMU.
The time method is simple and repeatable after calibration, while the IMU
method is more experimental because double integration accumulates drift.

\section{Wireless API and Web Interface}

The C6 robot creates its own Wi-Fi access point named
\texttt{CodeCell-C6-Drive}, so it does not require an external router. A phone
or laptop connects to the access point and opens the single-page controller.
The API is intentionally simple and can be tested directly from a browser.

\begin{center}
\begin{tabular}{@{}ll@{}}
\toprule
Endpoint & Function \\
\midrule
\texttt{/move} & Manual movement command \\
\texttt{/target} & Distance and angle command \\
\texttt{/status} & Battery, heading, distance, motor, and PID data \\
\texttt{/estop} & Immediate emergency stop \\
\texttt{/clear\_estop} & Clear the emergency stop state \\
\bottomrule
\end{tabular}
\end{center}

The HTML control app contains a touch-friendly virtual joystick, target-command
inputs, emergency buttons, and a live JSON status panel. It sends joystick
commands every 100 ms and polls status once per second. Releasing the joystick
automatically recentres it and sends a zero-speed command.

\section{Future ESP32-S3 Camera Subsystem}

The ESP32-S3 project was developed as a possible future camera subsystem. At
the current stage it is not installed on the robot and is not part of the
active drive platform. The firmware nevertheless configures the camera pins,
initializes the sensor in JPEG mode, starts a Wi-Fi access point named
\texttt{ESP32\_S3}, and serves a small camera web page. The camera firmware
exposes:

\begin{itemize}
  \item \texttt{/stream} on port 81 for an MJPEG live stream,
  \item \texttt{/capture} on port 80 for a single JPEG snapshot,
  \item \texttt{/control?cmd=...} for stream and camera power commands,
  \item \texttt{/control?cmd=status} for camera/stream state in JSON.
\end{itemize}

The camera can be turned off in software with \texttt{camera\_off}, which stops
the stream and deinitializes the camera driver. It can then be reinitialized
with \texttt{camera\_on}. If added in the future, this S3 subsystem could be
used for monitoring the robot without tying camera management to the real-time
drive controller.

\section{Safety and Robustness}

Several safeguards were implemented. Emergency stop is global and overrides
normal control by forcing both motor outputs to zero. The command timeout stops
the robot if the browser stops sending fresh direct-control commands. Shared
task state is protected by mutexes. Motor output is separated from control
calculation, so even when the communication task is busy, the motor task still
runs regularly. The browser joystick also recentres on pointer release, which
reduces the chance of leaving a nonzero command active accidentally.

\section{Summary}

The project produced a compact differential-drive robot with a 3D-printed
mechanical base, integrated C6 motor-control electronics, IMU-based heading
stabilisation, manual joystick driving, target-based movement, a browser REST
interface, and emergency handling. The ESP32-S3 camera firmware is prepared as
a future extension for live video and snapshots through a separate Wi-Fi camera
service. Therefore, the current platform is already usable as a C6-based mobile
robot, while future work can add obstacle avoidance, odometry, multi-waypoint
navigation, and integration between camera feedback and motion control.

\end{document}
